% Part: propositional-logic
% Chapter: propositional-logic 
% Section: soundness

\documentclass[../../../include/open-logic-section]{subfiles}

\begin{document}

\olfileid{pl}{prp}{snd}

\olsection{Kasahihan Logika Proposisional}

\begin{thm}[Kasahihan]
\ollabel{thm:soundness}
Yen $\Gamma\Proves!A$, mula $\Gamma \Entails !A$.
\end{thm}

\begin{proof}
Kanthi induksi marang carane teorema dibangun. Yen $!A$ minangka aksioma,
mula $\Entails !A$, mula mesthi $\Gamma \Entails!A$.
Semono uga, yen $!A \in \Gamma$, mula $\Gamma \Entails!A$.
Kanggo langkah induktif, upamane $!A$ diolehake kanthi \emph{modus ponens}
saka $!C$ lan $!C\lif!A$; uga anggepen minangka hipotesis induktif
yen teorema iki lumaku kanggo $!C\lif !A$ lan $!C$.
Butir katelu ing \olref[sem]{prop:semanticalfacts} menehi asil kang dikarepake.
\end{proof}

\begin{cor}
Yen $\Gamma$ bisa disembadani, mula $\Gamma$ konsisten.
Dadi, logika proposisional konsisten.
\end{cor}

\end{document}
